\chapter{Discussion}
\label{sec:discussion}
% Compare your results with the state-of-the-art and reflect upon the results and limitations of the study. You can already hint at future work to which you come back in the conclusion section.
%Write your discussion here. Do not forget to use sub-sections. Normally, the discussion starts with comparing your results to other studies as precisely as possible. The limitations should be reflected upon in terms such as reproducibility,  scalability,  generalizability,  reliability  and  validity. It is also important to mention ethical concerns.

The central finding is that automated server synthesis is a present capability, not a future prospect --- but its reliability is mediated by an underappreciated factor: whether the synthesis model's training-time representation of the target API is accurate and complete. This single axis explains the majority of variance in agent task performance across models, APIs, and documentation conditions. This chapter interprets the empirical findings in light of the research questions and prior work. Section~\ref{sec:discussion-synthesis} discusses the general feasibility of synthesis; Section~\ref{sec:discussion-gravity} examines the parametric gravity mechanism and its connection to the mechanistic probing results in Chapter~\ref{sec:mechanistic}; Section~\ref{sec:discussion-proxy} discusses SYNTH as a proxy metric; Section~\ref{sec:discussion-guidelines} derives practical guidelines for automated tool creation; and Section~\ref{sec:discussion-priorwork} positions the findings relative to prior work. Finally, key limitations are addressed in Section~\ref{sec:discussion-limitations}.

\section{Synthesis as a General Capability}
\label{sec:discussion-synthesis}

The central empirical finding is that LLM-based tool synthesis is broadly feasible: eleven of fourteen APIs achieve SYNTH $\geq 75\%$ in the docs condition, and the frontier-tier median agent pass rate of 83\% demonstrates that synthesized servers are not merely structurally valid but functionally useful. This is a non-trivial result. REST API integration involves strict interface contracts, multi-step authentication flows, and domain-specific
parameter conventions --- constraints that probabilistic generation does not inherently respect. That synthesis models navigate these requirements reliably across domains as diverse as e-commerce, communication, and financial data suggests the capability is not narrow or brittle. A five-run variance study on three APIs (Appendix~\ref{sec:apx:variance}) confirms that the no-docs condition is the most stable ($\sigma \approx 4$--$9$\,pp PASS across APIs), while the docs and mutated conditions show higher variance (up to $\sigma \approx 17$\,pp), primarily due to server size variation and degeneracy in complex APIs such as GitHub.

The two failure cases --- Jira and eBay Sell --- share a structural characteristic: extremely large documentation corpora with hundreds of endpoints distributed across deeply nested resource hierarchies. This points to a \emph{coverage scaling problem} rather than a fundamental inability to synthesize: the model understands the API but cannot attend to all endpoints
within a single synthesis context (see Appendix~\ref{sec:apx:complexity} for the relationship between documentation size and PASS). This is consistent with findings on positional bias in long contexts~\citep{liu2023lostmiddle,disentanglerecall2025},
and suggests that decomposition strategies --- synthesizing one API resource group at a time and assembling the results --- could close this gap without requiring larger context windows.

The comparison with human-authored reference servers on GitHub and Zulip sharpens this conclusion. Synthesized servers match the GitHub community server in the docs condition and substantially exceed it in the no-docs condition, while exceeding the Zulip reference in both conditions. The human references fail consistently on the same pattern: they cover the core, frequently-used API surface but omit less common but task-relevant endpoints (webhooks and Actions dispatch for GitHub; message edit and DM search for Zulip). A synthesis agent, by contrast, reads the full documentation and attempts comprehensive coverage --- and because it is not resource-constrained in the way a human contributor might be, this breadth comes without a corresponding quality trade-off.

\section{Parametric Gravity: Mechanism and Consequence}
\label{sec:discussion-gravity}

The mutation adherence results reveal a sharp split across APIs.
High-familiarity APIs (GitHub, Stripe, Shopify, Zulip, Jira) exhibit near-zero HTTP-layer adherence: the synthesis model ignores parameter mutations entirely and produces servers that reflect its training-time knowledge of the API. Low-familiarity APIs (Mastodon, Alphavantage) show near-complete adherence. This binary pattern is consistent with the training-frequency hypothesis of \citet{mallen2023knowns}: entities well-represented in pre-training data are retrieved from weights rather than from context.

The consequence for downstream performance is, however, more nuanced than a simple ``parametric reliance hurts'' narrative. For GitHub, zero adherence coincides with the highest PASS gains in the no-docs condition (+21\,pp for GPT-5.2, a single-model point estimate): the model writes from memory, and its memory is accurate. For Stripe, zero adherence coincides with a \emph{docs advantage} (+10\,pp for GPT-5.2): the model also writes from memory, but documentation provides additional signal that improves coverage. The mechanism that determines which outcome obtains is whether the synthesis model's training-time representation of the API is complete enough to be relied upon. API familiarity and API complexity interact here: GitHub is both highly familiar and structurally stable (a mature, slowly-evolving API), whereas Stripe is familiar but larger and has evolved significantly since the likely training cutoff of the primary models --- a scenario identified by \citet{lagging2025} as particularly susceptible to silent faithfulness failures. It is also worth noting that LLMs likely have strong priors for REST API patterns in general (HTTP verbs, JSON payloads, pagination conventions), even for less-exposed APIs, providing a structural scaffold that our mutation design does not directly probe.

This interpretation has a practical implication: whether to include documentation in the synthesis context is not a universal design choice but an API-specific one, and the right decision correlates with how confidently the synthesis model can write from memory. A pre-synthesis probe --- for example, asking the model to generate a list of API endpoints before providing documentation --- could serve as a low-cost signal for this confidence.

The mechanistic basis for the parametric gravity split is examined in Chapter~\ref{sec:mechanistic}. The logit-lens probing study confirms the binary pattern at the representational level: GitHub, Stripe, and eBay Commerce show 100\% parametric memory (the training-time attractor dominates within the first few transformer layers), while Mastodon, Confluence, and Jira show 100\% context-following. This early-to-mid network phenomenon aligns with feed-forward fact recall accounts~\citep{meng2022rome,factualrecall2025} and suggests that targeted context-faithful fine-tuning at transition layers could reduce parametric gravity without sacrificing synthesis capability. The finding also has implications for model choice: instruction fine-tuning and RLHF alter the geometry of the residual stream~\citep{posttrainingreshapes2025}, which may help explain the wide spread in adherence rates observed in Chapter~\ref{sec:results} --- GPT-5.4, for instance, reaches near-complete adherence on Confluence and eBay Commerce, while Gemini~3.1~Pro shows near-zero adherence across the benchmark, despite being large-scale models with comparable pre-training exposure (Gemini~3.1~Pro is classified as mid-tier in our benchmark based on synthesis reliability rather than raw model scale).

\section{Synthesis Quality as a Proxy Metric}
\label{sec:discussion-proxy}

The strong SYNTH--PASS correlation ($r = 0.937$, $p < 0.0001$, $n = 209$ across all synthesis models) has a direct practical implication: synthesis sufficiency --- assessed by an LLM judge examining the server's tool list without running the evaluation agent --- is a reliable predictor of downstream agent task performance. A judge sensitivity study across four models (Appendix~\ref{sec:apx:judge}) confirms that SYNTH scores are stable across judges: Gemini~3~Flash and GPT-4.1 agree within $\pm$12\,pp on the majority of conditions, validating the judge-based metric design. This matters because the full evaluation pipeline --- spinning up a sandboxed API environment, running the evaluation agent for up to 15 turns per task, calling the judge --- is expensive in both API quota and wall time. SYNTH can be assessed at synthesis time, before any agent evaluation occurs.

The two regimes visible in Figure~\ref{fig:scatter} sharpen when this metric is used diagnostically. For synthesis-limited cases (low SYNTH, low PASS), investment in improving synthesis quality --- better documentation pre-processing, larger context windows, iterative tool refinement --- will directly improve agent outcomes. For agent-limited cases (high SYNTH, lower PASS), the bottleneck is the evaluation agent's ability to compose available tools into correct multi-step plans, and improving synthesis will not help; the intervention should target the agent's reasoning or task decomposition.

One class of outlier is worth noting: Spoonacular in the docs condition achieves PASS=100\% despite SYNTH=75\%. This occurs because the evaluation agent compensates for minor coverage gaps using general culinary knowledge to reformulate queries through available tools. This agent-compensation effect is bounded --- it requires the agent to have its own parametric knowledge about the domain --- and cannot be relied upon for APIs in unfamiliar domains (compare Jira, where the agent has no comparable domain model to fall back on).

\section{Practical Design Guidelines}
\label{sec:discussion-guidelines}

The results point to several concrete recommendations for practitioners building full-server MCP synthesis pipelines --- systems that must produce complete, protocol-compliant tool servers rather than individual function stubs.

\paragraph{Documentation pre-processing matters for large APIs.}
For APIs with more than approximately 50 endpoints, raw documentation dumps produce coverage gaps: the synthesizer cannot attend to all endpoints in a single pass. Structured chunking by resource group, with per-chunk synthesis followed by assembly, is likely to produce higher SYNTH scores than end-to-end synthesis from a monolithic documentation corpus.

\paragraph{Omitting documentation is safe only for mature, well-known APIs.} For APIs with saturation-level representation in model training data --- GitHub, Mastodon, and major productivity platforms such as Notion --- the no-docs condition matches or exceeds the docs condition. This effect is partly about familiarity and partly about API maturity: the APIs where no-docs succeeds tend to be stable, slowly-evolving interfaces with few breaking version changes, so the model's training-time memory remains accurate. Documentation can add noise by providing competing context that the model must adjudicate against its priors, particularly for long, partially-outdated documentation files. For APIs at intermediate familiarity or with more active versioning (Stripe, Zulip), documentation provides a net benefit, and omitting it degrades performance. Practitioners should consider a pre-synthesis familiarity check --- for instance, asking the model to list API endpoints before receiving documentation --- before deciding whether to include documentation in the synthesis context.

\paragraph{Grounding manifests provide a weak but useful quality signal.} Requiring the synthesizer to produce a \texttt{grounding.json} enables a fast post-synthesis check: verify that every tool is anchored to an existing documentation file and a syntactically valid endpoint string. This does not guarantee implementation correctness --- a model can cite the right document while misreading it --- but tools lacking any documentation anchor are a reliable signal of unconstrained parametric generation, and missing grounding coverage (tools with no entry) often correlates with structural synthesis failures.

\paragraph{SYNTH is a cheap pre-flight metric.}
Running an LLM judge over the tool list before committing to full agent evaluation identifies synthesis-limited cases early. A SYNTH threshold of around 75\% appears to be the boundary below which full evaluation is unlikely to produce acceptable PASS rates in our benchmark; APIs below this threshold benefit more from re-synthesis than from agent tuning.

\paragraph{Prefer models with strong context-faithfulness for evolving or versioned APIs.} When the target API is less familiar to the model, or when documentation has been deliberately structured to override training priors, synthesis models that exhibit higher context adherence --- GPT-5.4 in our comparison, which reaches near-complete adherence on several mid-familiarity APIs --- produce more reliable servers. The GPT-5.4 regression on Stripe (8\% vs.\ 72\% baseline) illustrates the risk of deploying a model with strong parametric priors on an API that has evolved significantly since its training cutoff.

\paragraph{Looking ahead.} The guidelines above are grounded in the current single-agent synthesis scaffold. Future system designs could push further: neuro-symbolic verification (type-checking schemas against an API registry before committing), multi-agent pipelines with planner/verifier/assembler separation, and adaptive documentation inclusion (gated by a pre-synthesis familiarity probe) are natural extensions that the benchmark motivates but does not evaluate. These directions are discussed in the conclusion.

\section{Comparison with Prior Work}
\label{sec:discussion-priorwork}

Direct numerical comparison with prior work is complicated by the absence of overlapping evaluation sets: ToolFactory~\citep{toolfactory2025} and Doc2Agent~\citep{doc2agent2025} each evaluate on their own API selections and do not report PASS rates on the same 14 APIs used here. Qualitative comparison is nonetheless informative.

ToolFactory targets OpenAPI-specified inputs and focuses on coverage of the formal schema; our benchmark uses raw Markdown documentation, a substantially harder input modality that better reflects real practitioner conditions. Where ToolFactory reports synthesis success in terms of schema coverage, we report end-to-end agent task completion, which additionally exercises tool naming, description quality, and implementation correctness under live API execution.

Doc2Agent is closest to our work in motivation: it demonstrates that callable tool implementations can be generated from documentation at scale across many APIs. However, Doc2Agent generates individual tool stubs, not complete deployable servers; our synthesis target is a full MCP server with authentication, multi-tool coordination, and protocol compliance --- a more demanding integration target that is evaluated end-to-end through live API execution. Our benchmark additionally adds the no-docs and mutated conditions that Doc2Agent does not evaluate, enabling the parametric knowledge analysis that is our primary empirical contribution, and our comparison against the GitHub and Zulip community servers provides evidence on whether synthesized servers can match or exceed human-authored alternatives.

WAPIIBench~\citep{wapiibench2025} catalogues LLM-generated API integration errors against live services and reports error taxonomies that partially overlap with our failure cause categories (\verdict{tool\_schema}, \verdict{tool\_implementation}). Their finding that authentication and parameter serialisation errors are the most frequent failure types reflects a different primary failure regime than ours: WAPIIBench evaluates individual API calls, where the tool exists but its implementation is incorrect. Our primary failure mode is \verdict{tool\_coverage} --- the required tool was never synthesized at all --- so that when a tool is present, our \verdict{tool\_implementation} and \verdict{tool\_schema} rates are relatively low. This difference illustrates how full-server synthesis introduces a distinct and earlier failure point not captured by single-call evaluation.

\section{Limitations}
\label{sec:discussion-limitations}

\paragraph{Generalizability.}
The 14-API evaluation set spans five domains and a wide range of API sizes, but is not exhaustive. All selected APIs are REST-based, return JSON, and have English-language documentation. APIs with non-REST interfaces (GraphQL, gRPC), streaming responses, complex webhook architectures, or non-English documentation are not represented. The degree to which the parametric gravity findings generalise to programming-language SDKs rather than HTTP APIs is also unknown: the mutation adherence analysis is specific to the HTTP parameter-name setting.

\paragraph{Reliability.}
PASS rates are determined by an LLM judge, which introduces variability. We use a fixed judge model (Gemini~3~Flash) with a structured fixed-format prompt to reduce inconsistency. A judge sensitivity study across four judge models is reported in Appendix~\ref{sec:apx:judge} (Table~\ref{tab:apx:judge}); Gemini~3~Flash agrees with GPT-4.1 within $\pm$12\,pp on the majority of conditions. Similarly, the probing analysis rests on a single open-weight model; results should not be taken as established fact about frontier-scale models.

\paragraph{Reproducibility.}
All synthesis runs are stochastic: the same synthesis model and documentation may produce a different server on a different run. We report single-run results per condition rather than averages over multiple runs, which means some variance in PASS and SYNTH is unaccounted for; a synthesis variance study across five repeated runs on three APIs is reported in Appendix~\ref{sec:apx:variance} (Table~\ref{tab:apx:var_summary}). API sandboxes may also change state between runs, introducing noise in write-task evaluations.

\paragraph{Ethical considerations.}
Automated tool synthesis lowers the barrier to creating API integrations, including potentially for APIs whose terms of service restrict automated access or data extraction. The benchmark deliberately uses official sandbox environments with test credentials, but production deployment of synthesis pipelines against APIs not designed for automated access raises terms-of-service (ToS) and privacy concerns that practitioners should assess case by case. There is also an energy cost to running large synthesis models at scale; the no-docs condition, which is substantially cheaper, may be preferable where API familiarity is high, reducing unnecessary compute.